% Options for packages loaded elsewhere
\PassOptionsToPackage{unicode}{hyperref}
\PassOptionsToPackage{hyphens}{url}
\documentclass[11pt,
]{article}

\usepackage[margin=1in]{geometry}
\usepackage{setspace}
\setstretch{1.1}
\usepackage{parskip}

% Strict Widow and Club Penalties
\widowpenalty=10000
\clubpenalty=10000
\displaywidowpenalty=10000
\interlinepenalty=500
\brokenpenalty=10000

% Section Numbering
\setcounter{secnumdepth}{3}
\setcounter{tocdepth}{3}

\usepackage{xcolor}
\usepackage{amsmath,amssymb}
 % remove section numbering
\usepackage{iftex}
\ifPDFTeX
  \usepackage[T1]{fontenc}
  \usepackage[utf8]{inputenc}
  \usepackage{textcomp} % provide euro and other symbols
\else % if luatex or xetex
  \usepackage{unicode-math} % this also loads fontspec
  \defaultfontfeatures{Scale=MatchLowercase}
  \defaultfontfeatures[\rmfamily]{Ligatures=TeX,Scale=1}
\fi
\usepackage{lmodern}
\ifPDFTeX\else
  % xetex/luatex font selection
\fi
% Use upquote if available, for straight quotes in verbatim environments
\IfFileExists{upquote.sty}{\usepackage{upquote}}{}
\IfFileExists{microtype.sty}{% use microtype if available
  \usepackage[]{microtype}
  \UseMicrotypeSet[protrusion]{basicmath} % disable protrusion for tt fonts
}{}
\makeatletter
\@ifundefined{KOMAClassName}{% if non-KOMA class
  \IfFileExists{parskip.sty}{%
    \usepackage{parskip}
  }{% else
    \setlength{\parindent}{0pt}
    \setlength{\parskip}{6pt plus 2pt minus 1pt}}
}{% if KOMA class
  \KOMAoptions{parskip=half}}
\makeatother
\usepackage{longtable,booktabs,array}
\usepackage{caption}
\captionsetup[table]{skip=6pt}
\newcounter{none} % for unnumbered tables
\usepackage{calc} % for calculating minipage widths
% Correct order of tables after \paragraph or \subparagraph
\usepackage{etoolbox}
\makeatletter
\patchcmd\longtable{\par}{\if@noskipsec\mbox{}\fi\par}{}{}
\makeatother
% Allow footnotes in longtable head/foot
\IfFileExists{footnotehyper.sty}{\usepackage{footnotehyper}}{\usepackage{footnote}}
\makesavenoteenv{longtable}
\setlength{\emergencystretch}{3em} % prevent overfull lines
\providecommand{\tightlist}{%
  \setlength{\itemsep}{0pt}\setlength{\parskip}{0pt}}
\usepackage{bookmark}
\IfFileExists{xurl.sty}{\usepackage{xurl}}{} % add URL line breaks if available
\urlstyle{same}
% fallback for those not using the hyperref driver hyperxmp:
\makeatletter
\@ifundefined{xmpquote}{\newcommand{\xmpquote}[1]{#1}}{}
\makeatother
\hypersetup{
  hidelinks,
  pdfcreator={LaTeX via pandoc}}





\begin{document}
\begin{titlepage}
    \vspace*{2cm}
    \begin{center}
        {\Huge \textbf{FreshLens}} \\[1cm]
        {\LARGE AI Powered Freshness Monitoring System for Small Scale Retailers} \\[1.5cm]
        {\Large \textbf{Test Plan Report}} \\[3cm]
        {\Large Group 21} \\[0.5cm]
        {\Large PID 5}
    \end{center}
    \vfill
    \thispagestyle{empty}
\end{titlepage}
\newpage


{
\setcounter{tocdepth}{3}
\tableofcontents\newpage
}
\section{Evaluation Mission and Test
Motivation}\label{evaluation-mission-and-test-motivation}

Hey everyone! This Master Test Plan report covers our overall approach
to testing the FreshLens system. Since FreshLens is our core SaaS
platform for grocery retailers (and our main project for CS3203), we
really needed to make sure it works as expected under real-world
conditions. The main mission here is to provide a reliable, easy-to-use,
and highly secure online app for tracking inventory scans and freshness.
To give our end-users the best experience, we realized we had to design
and test this system pretty thoroughly across the board.

Since this is a web and mobile application handling people's business
data, it needs to be spot on. Even a single bug or a small security
loophole in the database could be a disaster for a vendor's inventory.
Our user interfaces and the backend logic have to be solid to deliver
exactly what the vendors expect. Since our requirements explicitly
outline multi-tenant data isolation, we had to put together a
well-organized testing process to verify that the system actually does
what it's supposed to do.

Our system is built using a modern architecture---mostly FastAPI for the
backend, Next.js for the web admin, and Expo for the mobile side. So our
testing plan has to cover all these different moving parts.

The main objectives of our testing process are:

\begin{itemize}
\tightlist
\item
  \textbf{Finding bugs:} We want to catch as many bugs as possible and
  fix them before we actually deploy this to our test users. We're doing
  a special search specifically for data leaks to make sure vendor CVs
  (like their company details) and inventory are completely
  confidential.
\item
  \textbf{Fixing design problems:} We need to find important problems in
  our design and implementation that might pose a threat to the quality
  of the system. For example, if our YOLO models are too slow, that's a
  problem we need to correct before the final delivery.
\item
  \textbf{Risk mitigation:} We're identifying risks associated with the
  project implementation. Things like failing user acceptance because
  the mobile camera lags. We test and refactor these before they happen
  by getting feedback early on.
\item
  \textbf{Quality standards:} The system has to keep up with existing
  standards. To achieve this, we're doing performance testing along with
  user acceptance testing to check if we're hitting our goals.
\item
  \textbf{Requirements check:} The system simply has to meet its
  requirements. A proper testing process helps us verify that we are
  actually building the functional and non-functional requirements we
  agreed upon, rather than wasting time on useless features.
\item
  \textbf{Legal/Process mandates:} We have to make sure we aren't
  violating any data privacy rules, so we're testing the system's legal
  and security boundaries. Every violation needs to be fixed before
  deployment.
\end{itemize}

\subsection{Out of Scope}\label{out-of-scope}

To keep our testing efforts focused on code our team actually wrote, the
following areas are explicitly \textbf{out of scope} for this test plan:

\begin{itemize}
\tightlist
\item
  \textbf{Supabase Authentication Infrastructure:} We test our parsing
  of JWTs, but we do not test the internal uptime, password hashing
  algorithms, or email delivery systems of Supabase itself.
\item
  \textbf{R2 / Cloud Storage Provider Internals:} We mock storage
  interactions or assume the bucket is available. The internal
  durability and replication of Cloudflare R2 are not tested by us.
\item
  \textbf{GitHub Actions Engine:} We rely on GitHub Actions for our CI
  pipeline, but we assume the underlying runner infrastructure is
  stable.
\end{itemize}

\subsection{Dependencies}\label{dependencies}

Our testing relies heavily on the following external and internal
dependencies to function correctly:

\begin{itemize}
\tightlist
\item
  \textbf{Working Test Accounts:} We require pre-seeded, active test
  accounts mapped to specific roles (\texttt{vendor},
  \texttt{platform\_admin}) and tenants to perform access control
  testing.
\item
  \textbf{CI Pipeline Uptime:} Automated tests require GitHub Actions to
  remain online.
\item
  \textbf{External Services:} Supabase endpoints (for JWKS retrieval)
  must stay reachable during integration test runs.
\end{itemize}

\subsection{Assumptions}\label{assumptions}

We are taking the following conditions as given during this testing
phase:

\begin{itemize}
\tightlist
\item
  \textbf{Stable Requirements:} The core functional requirements for ML
  inference and tenant isolation are stable and won't undergo massive
  refactors late in the cycle.
\item
  \textbf{Finalized RLS Policies:} The Postgres Row Level Security logic
  implemented in our raw SQL migrations is considered structurally
  finalized.
\item
  \textbf{Destructible Test Data:} Any data residing in the test
  environment databases can be freely created, corrupted, and destroyed
  by automated scripts without affecting stakeholders.
\end{itemize}

\subsection{Constraints}\label{constraints}

The following limits impact the extent of our testing effort:

\begin{itemize}
\tightlist
\item
  \textbf{No Dedicated QA Team:} Testing is performed entirely by the
  development team (us).
\item
  \textbf{No Prod-Matching Environment:} Due to budget constraints, our
  load testing environments (local Docker stacks) do not have the same
  CPU/RAM specifications as a true production cloud cluster.
\item
  \textbf{Time-Boxed Schedule:} Comprehensive testing is strictly
  time-boxed by our university project deadlines (M4, M5, M6).
\end{itemize}

\section{Target Test Items}\label{target-test-items}

Rather than a simple list, we've broken down our target test items into
a matrix that highlights the specific component and the main risk
associated with it if testing fails.

{\def\LTcaptype{none} % do not increment counter
\begin{longtable}[]{@{}
  >{\raggedright\arraybackslash}p{(\linewidth - 4\tabcolsep) * \real{0.3333}}
  >{\raggedright\arraybackslash}p{(\linewidth - 4\tabcolsep) * \real{0.3333}}
  >{\raggedright\arraybackslash}p{(\linewidth - 4\tabcolsep) * \real{0.3333}}@{}}
\toprule\noalign{}
\begin{minipage}[b]{\linewidth}\raggedright
Target Component
\end{minipage} & \begin{minipage}[b]{\linewidth}\raggedright
Description
\end{minipage} & \begin{minipage}[b]{\linewidth}\raggedright
Main Risk if untested/failed
\end{minipage} \\
\midrule\noalign{}
\endhead
\bottomrule\noalign{}
\endlastfoot
\textbf{User Interfaces (Mobile \& Web)} & The Next.js admin dashboard
and React Native Expo scanning app. & The vendor cannot upload scans due
to poor network logic, or the admin sees broken analytics charts. \\
\textbf{Data Models \& Database (Postgres)} & The core relational schema
and RLS policies (Postgres 16). & \textbf{Catastrophic cross-tenant data
leakage.} A vendor sees another vendor's stock. \\
\textbf{Functions (FastAPI \& ML Worker)} & The backend endpoints
(\texttt{/api/v1/sales}) and the Celery ML tasks. & Race conditions lead
to double stock deduction, or the ML model returns a false positive on
rotten produce. \\
\textbf{Event / Background Pipeline} & The Celery and Redis queuing
system managing asynchronous jobs. & Tasks are dropped silently during a
worker node crash, resulting in lost inventory scans. \\
\textbf{Release / CI Pipeline} & The GitHub actions workflow itself that
packages the builds. & A bad, failing build is accidentally shipped to
the main branch or production. \\
\end{longtable}
}

\section{Test Layering and Approach}\label{test-layering-and-approach}

We utilize a layered approach to testing, ensuring that fast, cheap
tests run frequently, while slower, expensive tests run at critical
milestones.

\subsection{Test Layering Table}\label{test-layering-table}

{\def\LTcaptype{none} % do not increment counter
\begin{longtable}[]{@{}
  >{\raggedright\arraybackslash}p{(\linewidth - 4\tabcolsep) * \real{0.3333}}
  >{\raggedright\arraybackslash}p{(\linewidth - 4\tabcolsep) * \real{0.3333}}
  >{\raggedright\arraybackslash}p{(\linewidth - 4\tabcolsep) * \real{0.3333}}@{}}
\toprule\noalign{}
\begin{minipage}[b]{\linewidth}\raggedright
Test Layer
\end{minipage} & \begin{minipage}[b]{\linewidth}\raggedright
Focus Area
\end{minipage} & \begin{minipage}[b]{\linewidth}\raggedright
CI Execution Timing
\end{minipage} \\
\midrule\noalign{}
\endhead
\bottomrule\noalign{}
\endlastfoot
\textbf{Unit} & Individual functions, Pydantic schemas, UI utilities
(e.g., chunked storage parser). & Every commit and PR. \\
\textbf{Integration} & Database connections, FastAPI endpoint logic
interacting with the database. & Every PR before merge. \\
\textbf{API} & Contract testing (checking requests/responses against
OpenAPI spec). & Every PR before merge. \\
\textbf{Interface (UI)} & Component rendering and form state logic in
React/React Native. & Every PR before merge. \\
\textbf{Cross-device} & Mobile layout rendering on iOS vs Android
(Expo). & Before major releases/milestones. \\
\textbf{Non-functional} & Load testing, RLS isolation security checks,
performance profiling. & RLS on every PR; Load testing before
release. \\
\textbf{Manual} & Camera access workflows, exploratory testing, UX
evaluations. & Pre-deployment staging review. \\
\end{longtable}
}

\subsection{Two-Dimensional Coverage
Definition}\label{two-dimensional-coverage-definition}

{\def\LTcaptype{none} % do not increment counter
\begin{longtable}[]{@{}
  >{\raggedright\arraybackslash}p{(\linewidth - 6\tabcolsep) * \real{0.2500}}
  >{\raggedright\arraybackslash}p{(\linewidth - 6\tabcolsep) * \real{0.2500}}
  >{\raggedright\arraybackslash}p{(\linewidth - 6\tabcolsep) * \real{0.2500}}
  >{\raggedright\arraybackslash}p{(\linewidth - 6\tabcolsep) * \real{0.2500}}@{}}
\toprule\noalign{}
\begin{minipage}[b]{\linewidth}\raggedright
Coverage Dimension
\end{minipage} & \begin{minipage}[b]{\linewidth}\raggedright
Definition
\end{minipage} & \begin{minipage}[b]{\linewidth}\raggedright
Tracking Method
\end{minipage} & \begin{minipage}[b]{\linewidth}\raggedright
Caveat
\end{minipage} \\
\midrule\noalign{}
\endhead
\bottomrule\noalign{}
\endlastfoot
\textbf{Requirement Coverage} & Ensuring every business rule (like
idempotent deductions) is explicitly validated by a test case. &
Generated directly from the pytest suites matching function names to
requirements. & Cannot guarantee a requirement is implemented
\emph{efficiently}. \\
\textbf{Code Coverage} & The percentage of lines of code executed during
the automated test suites. & Reported per-service via
\texttt{pytest-cov} and Jest coverage tools. & \textbf{100\% coverage
does not guarantee correctness;} it only proves the line was
executed. \\
\end{longtable}
}

\section{Testing Techniques and
Types}\label{testing-techniques-and-types}

\subsubsection{Data and Database Integrity
Testing}\label{data-and-database-integrity-testing}

{\def\LTcaptype{none} % do not increment counter
\begin{longtable}[]{@{}
  >{\raggedright\arraybackslash}p{(\linewidth - 2\tabcolsep) * \real{0.5000}}
  >{\raggedright\arraybackslash}p{(\linewidth - 2\tabcolsep) * \real{0.5000}}@{}}
\toprule\noalign{}
\begin{minipage}[b]{\linewidth}\raggedright
Attribute
\end{minipage} & \begin{minipage}[b]{\linewidth}\raggedright
Details
\end{minipage} \\
\midrule\noalign{}
\endhead
\bottomrule\noalign{}
\endlastfoot
\textbf{Technique Objective:} & Our database creation is done using SQL
migrations. We need to test this automatically by throwing incorrect
data and unauthorized queries at it. \\
\textbf{Technique:} & Invoke every entity using test scripts. Check that
correct data is retrieved (e.g., Tenant A only sees Tenant A's data). \\
\textbf{Oracles:} & Test access by injecting data through code, checking
behavior, and asserting retrieved data isn't erroneous. \\
\textbf{Required Tools:} & PostgreSQL 16, Pytest framework, custom SQL
isolation scripts (\texttt{rls\_isolation.sql}). \\
\textbf{Success Criteria:} & Every test is guaranteed to produce a
reasonable outcome, proving RLS bounds are perfectly respected. \\
\end{longtable}
}

\subsubsection{Function Testing}\label{function-testing}

{\def\LTcaptype{none} % do not increment counter
\begin{longtable}[]{@{}
  >{\raggedright\arraybackslash}p{(\linewidth - 2\tabcolsep) * \real{0.5000}}
  >{\raggedright\arraybackslash}p{(\linewidth - 2\tabcolsep) * \real{0.5000}}@{}}
\toprule\noalign{}
\begin{minipage}[b]{\linewidth}\raggedright
Attribute
\end{minipage} & \begin{minipage}[b]{\linewidth}\raggedright
Details
\end{minipage} \\
\midrule\noalign{}
\endhead
\bottomrule\noalign{}
\endlastfoot
\textbf{Technique Objective:} & Exercise the target functionality,
including data entry and processing, to log the behavior and check
requirement fulfillment. \\
\textbf{Technique:} & Execute each use-case scenario's individual flow
using valid and invalid data via Pytest. Verify appropriate error
messages and business rules. \\
\textbf{Oracles:} & Create inputs, determine expected outputs based on
OpenAPI specs, and assert the actual HTTP outcome. \\
\textbf{Required Tools:} & Pytest framework, FastAPI TestClient. \\
\textbf{Success Criteria:} & All tests execute properly without errors.
201 Created and 409 Conflict logic works flawlessly. \\
\end{longtable}
}

\subsubsection{User Interface Testing}\label{user-interface-testing}

{\def\LTcaptype{none} % do not increment counter
\begin{longtable}[]{@{}
  >{\raggedright\arraybackslash}p{(\linewidth - 2\tabcolsep) * \real{0.5000}}
  >{\raggedright\arraybackslash}p{(\linewidth - 2\tabcolsep) * \real{0.5000}}@{}}
\toprule\noalign{}
\begin{minipage}[b]{\linewidth}\raggedright
Attribute
\end{minipage} & \begin{minipage}[b]{\linewidth}\raggedright
Details
\end{minipage} \\
\midrule\noalign{}
\endhead
\bottomrule\noalign{}
\endlastfoot
\textbf{Technique Objective:} & Observe rendering with minimum memory,
high UX, platform independence, and smooth form submission workflows. \\
\textbf{Technique:} & Check whether complex logic (like chunked storage
uploads) works properly under throttled network settings. \\
\textbf{Oracles:} & Automate testing for UI utilities
(\texttt{claims.test.ts}), while relying on manual walkthroughs for
complex gestures. \\
\textbf{Required Tools:} & Jest for React testing, ESLint, browser dev
tools. \\
\textbf{Success Criteria:} & All routes resolve successfully, zero
broken components, no memory leak warnings in console. \\
\end{longtable}
}

\subsubsection{Load and Performance
Testing}\label{load-and-performance-testing}

{\def\LTcaptype{none} % do not increment counter
\begin{longtable}[]{@{}
  >{\raggedright\arraybackslash}p{(\linewidth - 2\tabcolsep) * \real{0.5000}}
  >{\raggedright\arraybackslash}p{(\linewidth - 2\tabcolsep) * \real{0.5000}}@{}}
\toprule\noalign{}
\begin{minipage}[b]{\linewidth}\raggedright
Attribute
\end{minipage} & \begin{minipage}[b]{\linewidth}\raggedright
Details
\end{minipage} \\
\midrule\noalign{}
\endhead
\bottomrule\noalign{}
\endlastfoot
\textbf{Technique Objective:} & Exercise transactions under normal,
worst-case, and concurrent user workloads to observe system performance
degradation. \\
\textbf{Technique:} & Run loops in test scripts to simultaneously hit
the API, increasing transaction volume to simulate spiky traffic. \\
\textbf{Oracles:} & Code must be automated to apply the load, monitoring
CPU and database connection pool saturation. \\
\textbf{Required Tools:} & Custom async Python scripts and Pytest
concurrency fixtures. \\
\textbf{Success Criteria:} & The system gracefully degrades or queues
requests rather than dropping them or returning 500 errors. \\
\end{longtable}
}

\subsubsection{Security and Access Control
Testing}\label{security-and-access-control-testing}

{\def\LTcaptype{none} % do not increment counter
\begin{longtable}[]{@{}
  >{\raggedright\arraybackslash}p{(\linewidth - 2\tabcolsep) * \real{0.5000}}
  >{\raggedright\arraybackslash}p{(\linewidth - 2\tabcolsep) * \real{0.5000}}@{}}
\toprule\noalign{}
\begin{minipage}[b]{\linewidth}\raggedright
Attribute
\end{minipage} & \begin{minipage}[b]{\linewidth}\raggedright
Details
\end{minipage} \\
\midrule\noalign{}
\endhead
\bottomrule\noalign{}
\endlastfoot
\textbf{Technique Objective:} & Ensure an actor can access only
functions they have permissions for, and unauthenticated users are
entirely blocked. \\
\textbf{Technique:} & Create tests for each user type (\texttt{vendor},
\texttt{platform\_admin}). Verify forbidden messages when attempting
unauthorized access. \\
\textbf{Oracles:} & Fully automated testing using Pytest to assert HTTP
401 Unauthorized and 403 Forbidden response codes. \\
\textbf{Required Tools:} & Pytest, PyJWT for generating mock tokens. \\
\textbf{Success Criteria:} & System rigidly enforces privileges; access
denied messages are given strictly and correctly. \\
\end{longtable}
}

\section{Evidence, Reporting, and
Deliverables}\label{evidence-reporting-and-deliverables}

We have established strict mechanics for how test evidence is captured
and retained to ensure maximum transparency.

\subsection{Evidence Retention
Mechanics}\label{evidence-retention-mechanics}

\begin{itemize}
\tightlist
\item
  \textbf{Machine-Readable Result Files:} All test runs produce XML/JSON
  output files (like JUnit XML) that are archived as artifacts in GitHub
  Actions for historic trend analysis.
\item
  \textbf{Raw Request Logs:} During load testing, we do not simply
  retain the ``summary'' report. Raw HTTP request logs (including
  latency and status codes) are retained in storage to allow for
  deep-dive debugging of latency spikes.
\item
  \textbf{UI Test Failures:} (Future implementation) If a UI test fails,
  the CI pipeline will capture a screenshot/video of the headless
  browser state to immediately highlight the visual error.
\end{itemize}

\subsection{Explicit Rules on
Credentials}\label{explicit-rules-on-credentials}

\textbf{Strict Rule:} No production or staging credentials, API keys, or
raw passwords will ever be committed to the repository for testing
purposes.

\begin{itemize}
\tightlist
\item
  All test suites must rely on mock generation, factory fakes, or rely
  on \texttt{.env.test} templates with blank fields. Real secrets are
  injected only at runtime by the CI/CD secure vault.
\end{itemize}

\subsection{Calling Out Weakest Areas}\label{calling-out-weakest-areas}

Rather than hiding behind a blended coverage percentage (e.g., ``The
project has 85\% coverage''), our test evaluation reports explicitly
name the ``weakest areas'' that lack coverage. This ensures we are
always aware of technical debt (for instance, calling out that mobile
offline caching currently lacks automated coverage).

\subsection{Test Evaluation Summaries}\label{test-evaluation-summaries}

\begin{itemize}
\tightlist
\item
  \textbf{Test Logs:} These are the detailed logs that give the outcome
  of every test written for the functionalities. If even a single test
  log gives an output as a failure in our GitHub Actions, that
  functionality is recorrected and all tests are run again.
\item
  \textbf{Code Inspection Tool:} Code inspection details are gotten
  after running ESLint and TypeScript checkers in the IDE. This notifies
  us about spelling mistakes, empty tags, and unused items.
\item
  \textbf{Database Performance Statistics:} This is a very good method
  to check database performance. It checks memory usage and size.
\end{itemize}

\section{Detailed Test Cases}\label{detailed-test-cases}

To provide a much more thorough explanation of our testing workflow, we
have detailed our most critical test cases below. Rather than a brief
overview, these describe exactly why and how we test each functional
component of the FreshLens system to ensure it holds up to real-world
vendor usage.

\subsection{Backend API Test Case: Scan Upload Workflow
(API-01)}\label{backend-api-test-case-scan-upload-workflow-api-01}

\begin{itemize}
\tightlist
\item
  \textbf{Motivation:} Vendors operating in grocery stores often have
  terrible Wi-Fi or cellular connections. If they upload a large image
  of produce, it cannot time out and crash the application.
\item
  \textbf{Pre-conditions:} The mobile app is authenticated with a valid
  \texttt{vendor} JWT token. The server is online and connected to the
  Redis broker.
\item
  \textbf{Execution Steps:}

  \begin{enumerate}
  \def\labelenumi{\arabic{enumi}.}
  \tightlist
  \item
    The vendor opens the camera within the Expo app.
  \item
    The vendor takes a high-resolution photo of a banana.
  \item
    The photo is chunked locally and transmitted via HTTP POST to
    \texttt{/api/v1/scans}.
  \end{enumerate}
\item
  \textbf{Expected System Behavior:} The FastAPI backend should
  immediately capture the image byte chunks, save them to the R2 bucket,
  and spawn an asynchronous Celery task. Critically, it must immediately
  return a \texttt{202\ Accepted} status code to the mobile app so the
  vendor isn't left waiting for the ML inference to finish.
\item
  \textbf{Status:} \textbf{PASS.} Verified via \texttt{test\_scans.py}.
\end{itemize}

\subsection{Backend API Test Case: Sales Idempotency
(API-02)}\label{backend-api-test-case-sales-idempotency-api-02}

\begin{itemize}
\tightlist
\item
  \textbf{Motivation:} If a vendor hits the ``Submit Sale'' button twice
  by accident, or if the network stutters and resends the request, we
  absolutely cannot deduct stock twice. This would ruin the vendor's
  inventory records.
\item
  \textbf{Pre-conditions:} The database has an existing inventory batch
  of tomatoes with a quantity greater than zero.
\item
  \textbf{Execution Steps:}

  \begin{enumerate}
  \def\labelenumi{\arabic{enumi}.}
  \tightlist
  \item
    A script generates a unique \texttt{Idempotency-Key} header.
  \item
    The script fires two completely identical POST requests
    simultaneously to \texttt{/api/v1/sales}.
  \end{enumerate}
\item
  \textbf{Expected System Behavior:} The backend must rely on a database
  transaction lock or caching layer. The very first request processed
  should succeed and return \texttt{201\ Created}, deducting the stock
  exactly once. The second concurrent request should hit the idempotency
  lock and gracefully return a \texttt{409\ Conflict}, without modifying
  the database.
\item
  \textbf{Status:} \textbf{PASS.} Verified via \texttt{test\_sales.py}.
\end{itemize}

\subsection{Database Security Test Case: Tenant Boundary Isolation
(SEC-01)}\label{database-security-test-case-tenant-boundary-isolation-sec-01}

\begin{itemize}
\tightlist
\item
  \textbf{Motivation:} Since this is a multi-tenant SaaS application,
  multiple different grocery retailers are sharing the same database
  tables. If Tenant A can somehow view or edit Tenant B's sales data,
  the entire platform's security is compromised.
\item
  \textbf{Pre-conditions:} The Postgres 16 database is seeded with data
  for both ``Tenant A'' and ``Tenant B''. Row Level Security (RLS) is
  active.
\item
  \textbf{Execution Steps:}

  \begin{enumerate}
  \def\labelenumi{\arabic{enumi}.}
  \tightlist
  \item
    An automated SQL script logs into the database using the restricted
    \texttt{freshlens\_api\_runtime} role.
  \item
    The script sets the local transaction context to Tenant A.
  \item
    The script attempts to run
    \texttt{SELECT\ *\ FROM\ scans\ WHERE\ tenant\_id\ =\ \textquotesingle{}Tenant\_B\_ID\textquotesingle{}}.
  \end{enumerate}
\item
  \textbf{Expected System Behavior:} Instead of crashing, the Postgres
  RLS engine silently filters out the unauthorized rows, returning an
  empty list (0 rows) to Tenant A. The application backend is entirely
  shielded from accidentally leaking data.
\item
  \textbf{Status:} \textbf{PASS.} Verified via
  \texttt{infra/db/tests/rls\_isolation.sql}.
\end{itemize}

\subsection{ML Pipeline Test Case: Identity Confidence Threshold
(ML-01)}\label{ml-pipeline-test-case-identity-confidence-threshold-ml-01}

\begin{itemize}
\tightlist
\item
  \textbf{Motivation:} The YOLO object detection model is incredibly
  fast but can sometimes be overconfident when analyzing blurry images
  or items it has never seen before. We cannot let it misclassify a
  blurry eggplant as a cucumber and ruin inventory counts.
\item
  \textbf{Pre-conditions:} The \texttt{YOLO26n-cls} model is loaded into
  the Celery worker memory.
\item
  \textbf{Execution Steps:}

  \begin{enumerate}
  \def\labelenumi{\arabic{enumi}.}
  \tightlist
  \item
    A known ambiguous or blurry image is passed directly to the
    inference pipeline via the testing framework.
  \end{enumerate}
\item
  \textbf{Expected System Behavior:} The model evaluates the image. Even
  if its top prediction is ``cucumber'', it must check the confidence
  score. Because the confidence score is strictly enforced to be
  \texttt{\textgreater{}=\ 0.75} for identity, the script must intercept
  a lower score (e.g., 0.60) and override the prediction, safely
  returning an \texttt{unknown} label instead of logging a false
  positive.
\item
  \textbf{Status:} \textbf{PASS.} Verified via
  \texttt{test\_identity.py}.
\end{itemize}

\subsection{UI Test Case: Admin Dashboard Analytics Rendering
(UI-01)}\label{ui-test-case-admin-dashboard-analytics-rendering-ui-01}

\begin{itemize}
\tightlist
\item
  \textbf{Motivation:} The platform administrator needs to view
  aggregate data. If the TypeScript data mapping logic fails when
  receiving an unexpected payload from the API, the entire React screen
  will crash (White Screen of Death).
\item
  \textbf{Pre-conditions:} The user is logged in via Supabase with the
  \texttt{app\_role} explicitly set to \texttt{platform\_admin}.
\item
  \textbf{Execution Steps:}

  \begin{enumerate}
  \def\labelenumi{\arabic{enumi}.}
  \tightlist
  \item
    The admin navigates to the core analytics dashboard URL.
  \item
    The page fires a fetch request for aggregate mapping data.
  \end{enumerate}
\item
  \textbf{Expected System Behavior:} The Next.js frontend utilizes
  \texttt{admin-map.test.ts} logic to parse the arrays. It should map
  the data without throwing any TypeScript \texttt{undefined} errors and
  cleanly render the visualization chart, displaying aggregated
  cross-tenant metrics flawlessly.
\item
  \textbf{Status:} \textbf{PASS.} Verified via Jest UI suites.
\end{itemize}

\section{Appendix: Implemented Test
Inventory}\label{appendix-implemented-test-inventory}

Here is the exhaustive inventory of all test suites currently
implemented in our project. We wrote these to make sure our system is as
robust as possible.

\subsection{\texorpdfstring{Backend API Tests
(\texttt{apps/api/tests/})}{Backend API Tests (apps/api/tests/)}}\label{backend-api-tests-appsapitests}

\begin{itemize}
\tightlist
\item
  \texttt{test\_admin.py}
\item
  \texttt{test\_scans.py}
\item
  \texttt{test\_auth.py}
\item
  \texttt{test\_storage.py}
\item
  \texttt{test\_authorization.py}
\item
  \texttt{test\_config.py}
\item
  \texttt{test\_health.py}
\item
  \texttt{test\_sales.py}
\item
  \texttt{test\_tenant\_context.py}
\end{itemize}

\subsection{\texorpdfstring{Machine Learning Worker Tests
(\texttt{packages/ml/tests/})}{Machine Learning Worker Tests (packages/ml/tests/)}}\label{machine-learning-worker-tests-packagesmltests}

\begin{itemize}
\tightlist
\item
  \texttt{test\_prepare\_identity\_dataset.py}
\item
  \texttt{test\_stub\_classifier.py}
\item
  \texttt{test\_evaluate\_freshness.py}
\item
  \texttt{test\_inventory\_batch.py}
\item
  \texttt{test\_imagenet\_produce.py}
\item
  \texttt{test\_prepare\_freshness\_dataset.py}
\item
  \texttt{test\_identity.py}
\item
  \texttt{test\_dataset\_snapstock.py}
\end{itemize}

\subsection{\texorpdfstring{Frontend \& Mobile Logic Tests
(\texttt{apps/web/} \&
\texttt{apps/mobile/})}{Frontend \& Mobile Logic Tests (apps/web/ \& apps/mobile/)}}\label{frontend-mobile-logic-tests-appsweb-appsmobile}

\begin{itemize}
\tightlist
\item
  \texttt{claims.test.ts} (Web)
\item
  \texttt{admin-map.test.ts} (Web)
\item
  \texttt{claims.test.ts} (Mobile)
\item
  \texttt{chunked-storage.test.ts} (Mobile)
\end{itemize}

\subsection{\texorpdfstring{Database Security Tests
(\texttt{infra/db/tests/})}{Database Security Tests (infra/db/tests/)}}\label{database-security-tests-infradbtests}

\begin{itemize}
\tightlist
\item
  \texttt{rls\_isolation.sql}
\item
  \texttt{runtime\_role.sql}
\end{itemize}

\newpage\section{References}\label{references}

\begin{itemize}
\tightlist
\item
  FreshLens System Architecture Document
  (\texttt{docs/design/FreshLens-SAD.md})
\item
  FreshLens API Contract (\texttt{docs/api/v1/openapi.yaml})
\item
  Pytest Testing Framework available at https://docs.pytest.org/
\item
  Jest Testing Framework available at https://jestjs.io/
\item
  Supabase Authentication Architecture available at
  https://supabase.com/docs/guides/auth
\end{itemize}

\end{document}
